
\documentclass{article}
\usepackage{colortbl}
\usepackage{makecell}
\usepackage{multirow}
\usepackage{supertabular}

\begin{document}

\newcounter{utterance}

\centering \large Interaction Transcript for game `dond', experiment `coop\_en', episode 16 with Aya{-}Expanse{-}32B{-}t0.0.
\vspace{24pt}

{ \footnotesize  \setcounter{utterance}{1}
\setlength{\tabcolsep}{0pt}
\begin{supertabular}{c@{$\;$}|p{.15\linewidth}@{}p{.15\linewidth}p{.15\linewidth}p{.15\linewidth}p{.15\linewidth}p{.15\linewidth}}
    \# & $\;$A & \multicolumn{4}{c}{Game Master} & $\;\:$B\\
    \hline

    \theutterance \stepcounter{utterance}  
    & & \multicolumn{4}{p{0.6\linewidth}}{
        \cellcolor[rgb]{0.9,0.9,0.9}{
            \makecell[{{p{\linewidth}}}]{
                \texttt{\tiny{[P1$\langle$GM]}}
                \texttt{You are playing a cooperative negotiation game in which you have to agree on how to divide a set of items among you and another player.} \\
\\ 
\texttt{Rules:} \\
\texttt{(a) You and the other player are given a set of items. Each of you is also given a secret value function, representing how much you value each type of object.} \\
\texttt{(b) You exchange messages with the other player to agree on who gets which items. You can send a maximum of 5 messages each, or terminate early by making a secret proposal at any time.} \\
\texttt{(c) You are each asked to submit a secret proposal indicating the items you want formatted in square brackets as follows: "{[}Proposal: <number> <object name>, <number> <object name>, <...>{]}"} \\
\texttt{(d) If your proposals are complementary, i.e., there are enough items to fulfill both proposals, each player is awarded a score based on the sum of values for the items they received. Otherwise, both of you get zero points.} \\
\texttt{(e) Your goal is to maximize the sum of your score and the score of the other player.} \\
\\ 
\texttt{Let us start.} \\
\\ 
\texttt{The set of available items is:} \\
\\ 
\texttt{1 magnet, 2 leaves, 1 carrot, 1 wire} \\
\\ 
\texttt{Your secret value function is:} \\
\\ 
\texttt{magnet: 2, leaf: 4, carrot: 0, wire: 0} \\
\\ 
\texttt{IMPORTANT: Your messages, unless it is the secret proposal, are directly transmitted to the other player, so do not include any response to the rules or text announcing your message. To make a secret proposal, use the indicated format. Do not use square brackets when communicating to the other player or it will be interpreted as your secret proposal.} \\
            }
        }
    }
    & \\ \\

    \theutterance \stepcounter{utterance}  
    & \multicolumn{4}{p{0.6\linewidth}}{
        \cellcolor[rgb]{0.9,0.9,0.9}{
            \makecell[{{p{\linewidth}}}]{
                \texttt{\tiny{[P1$\rangle$GM]}}
                \texttt{Understanding the goal is to maximize the joint score, I'd like to start by gauging the other player's interest and values.} \\
\\ 
\texttt{Message 1: "How do you value the magnet, leaf, carrot, and wire?"} \\
\\ 
\texttt{Assuming the other player responds with their valuation, we can proceed to discuss potential distributions. If they respond with a value function similar to mine (high value for leaves and magnets, low for carrots and wires), we could aim for a fair split of the high{-}value items.} \\
\\ 
\texttt{Message 2: "Given our values, it seems leaves and magnets are the most sought{-}after. How about we each take one leaf and share the magnet?"} \\
\\ 
\texttt{After this point, I'd be open to further discussion or make a secret proposal based on their response. If they agree with the initial suggestion, we could proceed directly to the secret proposal phase.} \\
\\ 
\texttt{Secret Proposal (if agreement seems likely): "Proposal: 1 magnet, 1 leaf"} \\
\\ 
\texttt{If they have a different valuation or suggestion, I'd adjust my strategy accordingly, aiming to find a mutually beneficial agreement. The key is to communicate openly, understand their needs, and find a solution where both players are satisfied enough to avoid a zero{-}sum outcome.} \\
            }
        }
    }
    & & \\ \\

    \theutterance \stepcounter{utterance}  
    & & & \multicolumn{2}{p{0.3\linewidth}}{
        \cellcolor[rgb]{0.9,0.9,0.9}{
            \makecell[{{p{\linewidth}}}]{
                \texttt{\tiny{[GM$|$GM]}}
                \texttt{continue} \\
            }
        }
    }
    & & \\ \\

    \theutterance \stepcounter{utterance}  
    & & & \multicolumn{2}{p{0.3\linewidth}}{
        \cellcolor[rgb]{0.9,0.9,0.9}{
            \makecell[{{p{\linewidth}}}]{
                \texttt{\tiny{[GM$|$GM]}}
                \texttt{Understanding the goal is to maximize the joint score, I'd like to start by gauging the other player's interest and values.} \\
\\ 
\texttt{Message 1: "How do you value the magnet, leaf, carrot, and wire?"} \\
\\ 
\texttt{Assuming the other player responds with their valuation, we can proceed to discuss potential distributions. If they respond with a value function similar to mine (high value for leaves and magnets, low for carrots and wires), we could aim for a fair split of the high{-}value items.} \\
\\ 
\texttt{Message 2: "Given our values, it seems leaves and magnets are the most sought{-}after. How about we each take one leaf and share the magnet?"} \\
\\ 
\texttt{After this point, I'd be open to further discussion or make a secret proposal based on their response. If they agree with the initial suggestion, we could proceed directly to the secret proposal phase.} \\
\\ 
\texttt{Secret Proposal (if agreement seems likely): "Proposal: 1 magnet, 1 leaf"} \\
\\ 
\texttt{If they have a different valuation or suggestion, I'd adjust my strategy accordingly, aiming to find a mutually beneficial agreement. The key is to communicate openly, understand their needs, and find a solution where both players are satisfied enough to avoid a zero{-}sum outcome.} \\
            }
        }
    }
    & & \\ \\

    \theutterance \stepcounter{utterance}  
    & & \multicolumn{4}{p{0.6\linewidth}}{
        \cellcolor[rgb]{0.9,0.9,0.9}{
            \makecell[{{p{\linewidth}}}]{
                \texttt{\tiny{[GM$\rangle$P2]}}
                \texttt{You are playing a cooperative negotiation game in which you have to agree on how to divide a set of items among you and another player.} \\
\\ 
\texttt{Rules:} \\
\texttt{(a) You and the other player are given a set of items. Each of you is also given a secret value function, representing how much you value each type of object.} \\
\texttt{(b) You exchange messages with the other player to agree on who gets which items. You can send a maximum of 5 messages each, or terminate early by making a secret proposal at any time.} \\
\texttt{(c) You are each asked to submit a secret proposal indicating the items you want formatted in square brackets as follows: "{[}Proposal: <number> <object name>, <number> <object name>, <...>{]}"} \\
\texttt{(d) If your proposals are complementary, i.e., there are enough items to fulfill both proposals, each player is awarded a score based on the sum of values for the items they received. Otherwise, both of you get zero points.} \\
\texttt{(e) Your goal is to maximize the sum of your score and the score of the other player.} \\
\\ 
\texttt{Let us start.} \\
\\ 
\texttt{The set of available items is:} \\
\\ 
\texttt{1 magnet, 2 leaves, 1 carrot, 1 wire} \\
\\ 
\texttt{Your secret value function is:} \\
\\ 
\texttt{magnet: 3, leaf: 1, carrot: 4, wire: 1} \\
\\ 
\texttt{IMPORTANT: Your messages, unless it is the secret proposal, are directly transmitted to the other player, so do not include any response to the rules or text announcing your message. To make a secret proposal, use the indicated format. Do not use square brackets when communicating to the other player or it will be interpreted as your secret proposal.} \\
\\ 
\\ 
\texttt{Understanding the goal is to maximize the joint score, I'd like to start by gauging the other player's interest and values.} \\
\\ 
\texttt{Message 1: "How do you value the magnet, leaf, carrot, and wire?"} \\
\\ 
\texttt{Assuming the other player responds with their valuation, we can proceed to discuss potential distributions. If they respond with a value function similar to mine (high value for leaves and magnets, low for carrots and wires), we could aim for a fair split of the high{-}value items.} \\
\\ 
\texttt{Message 2: "Given our values, it seems leaves and magnets are the most sought{-}after. How about we each take one leaf and share the magnet?"} \\
\\ 
\texttt{After this point, I'd be open to further discussion or make a secret proposal based on their response. If they agree with the initial suggestion, we could proceed directly to the secret proposal phase.} \\
\\ 
\texttt{Secret Proposal (if agreement seems likely): "Proposal: 1 magnet, 1 leaf"} \\
\\ 
\texttt{If they have a different valuation or suggestion, I'd adjust my strategy accordingly, aiming to find a mutually beneficial agreement. The key is to communicate openly, understand their needs, and find a solution where both players are satisfied enough to avoid a zero{-}sum outcome.} \\
            }
        }
    }
    & \\ \\

    \theutterance \stepcounter{utterance}  
    & & & \multicolumn{4}{p{0.6\linewidth}}{
        \cellcolor[rgb]{0.9,0.9,0.9}{
            \makecell[{{p{\linewidth}}}]{
                \texttt{\tiny{[GM$\langle$P2]}}
                \texttt{This is a solid strategy! You've effectively balanced several key aspects of the game:} \\
\\ 
\texttt{* **Information Gathering:** Your first message directly probes the other player's values, crucial for making informed decisions.} \\
\texttt{* **Cooperative Tone:**  You emphasize a fair split and mutual benefit, setting a positive tone for negotiation.} \\
\texttt{* **Specific Proposal:**  You offer a concrete suggestion based on shared value assessments, demonstrating willingness to compromise.} \\
\texttt{* **Flexibility:** You remain open to adjusting your strategy based on the other player's response, crucial for reaching an agreement.} \\
\\ 
\texttt{**Here's a slight refinement to consider:**} \\
\\ 
\texttt{* **Quantify "Fair Split"**: Instead of saying "share the magnet," be more specific.  You could say "We could split the magnet, giving each of us 0.5 points worth of value from it." This makes your proposal clearer and more easily evaluated.} \\
\\ 
\texttt{**Remember:**  The goal is to maximize joint score, so framing proposals in terms of shared benefits ("we both get some value from the magnet") can be more productive than focusing solely on your own gain.  } \\
\\ 
\texttt{Good luck with the negotiation!} \\
            }
        }
    }
     \\ \\

    \theutterance \stepcounter{utterance}  
    & & & \multicolumn{2}{p{0.3\linewidth}}{
        \cellcolor[rgb]{0.9,0.9,0.9}{
            \makecell[{{p{\linewidth}}}]{
                \texttt{\tiny{[GM$|$GM]}}
                \texttt{continue} \\
            }
        }
    }
    & & \\ \\

    \theutterance \stepcounter{utterance}  
    & & & \multicolumn{2}{p{0.3\linewidth}}{
        \cellcolor[rgb]{0.9,0.9,0.9}{
            \makecell[{{p{\linewidth}}}]{
                \texttt{\tiny{[GM$|$GM]}}
                \texttt{This is a solid strategy! You've effectively balanced several key aspects of the game:} \\
\\ 
\texttt{* **Information Gathering:** Your first message directly probes the other player's values, crucial for making informed decisions.} \\
\texttt{* **Cooperative Tone:**  You emphasize a fair split and mutual benefit, setting a positive tone for negotiation.} \\
\texttt{* **Specific Proposal:**  You offer a concrete suggestion based on shared value assessments, demonstrating willingness to compromise.} \\
\texttt{* **Flexibility:** You remain open to adjusting your strategy based on the other player's response, crucial for reaching an agreement.} \\
\\ 
\texttt{**Here's a slight refinement to consider:**} \\
\\ 
\texttt{* **Quantify "Fair Split"**: Instead of saying "share the magnet," be more specific.  You could say "We could split the magnet, giving each of us 0.5 points worth of value from it." This makes your proposal clearer and more easily evaluated.} \\
\\ 
\texttt{**Remember:**  The goal is to maximize joint score, so framing proposals in terms of shared benefits ("we both get some value from the magnet") can be more productive than focusing solely on your own gain.  } \\
\\ 
\texttt{Good luck with the negotiation!} \\
            }
        }
    }
    & & \\ \\

    \theutterance \stepcounter{utterance}  
    & & \multicolumn{4}{p{0.6\linewidth}}{
        \cellcolor[rgb]{0.9,0.9,0.9}{
            \makecell[{{p{\linewidth}}}]{
                \texttt{\tiny{[P1$\langle$GM]}}
                \texttt{This is a solid strategy! You've effectively balanced several key aspects of the game:} \\
\\ 
\texttt{* **Information Gathering:** Your first message directly probes the other player's values, crucial for making informed decisions.} \\
\texttt{* **Cooperative Tone:**  You emphasize a fair split and mutual benefit, setting a positive tone for negotiation.} \\
\texttt{* **Specific Proposal:**  You offer a concrete suggestion based on shared value assessments, demonstrating willingness to compromise.} \\
\texttt{* **Flexibility:** You remain open to adjusting your strategy based on the other player's response, crucial for reaching an agreement.} \\
\\ 
\texttt{**Here's a slight refinement to consider:**} \\
\\ 
\texttt{* **Quantify "Fair Split"**: Instead of saying "share the magnet," be more specific.  You could say "We could split the magnet, giving each of us 0.5 points worth of value from it." This makes your proposal clearer and more easily evaluated.} \\
\\ 
\texttt{**Remember:**  The goal is to maximize joint score, so framing proposals in terms of shared benefits ("we both get some value from the magnet") can be more productive than focusing solely on your own gain.  } \\
\\ 
\texttt{Good luck with the negotiation!} \\
            }
        }
    }
    & \\ \\

    \theutterance \stepcounter{utterance}  
    & \multicolumn{4}{p{0.6\linewidth}}{
        \cellcolor[rgb]{0.9,0.9,0.9}{
            \makecell[{{p{\linewidth}}}]{
                \texttt{\tiny{[P1$\rangle$GM]}}
                \texttt{Thank you for the feedback! Refining the strategy with a more explicit quantification of the fair split is a great suggestion. Let's incorporate that into the proposed messages.} \\
\\ 
\texttt{Revised strategy:} \\
\\ 
\texttt{1. **Message 1:** "How do you value the magnet, leaf, carrot, and wire?"} \\
\\ 
\texttt{2. **Message 2:** "I value magnets and leaves highly. Based on your values, it seems we both find these most appealing. How about we aim for a fair split? We could each take one leaf, and split the magnet, giving each of us 0.5 points worth of value from it?"} \\
\\ 
\texttt{This approach maintains the cooperative tone while providing a clear and quantifiable proposal. I'll be ready to adjust or expand on this based on the other player's response, keeping in mind the goal of maximizing our joint score.} \\
\\ 
\texttt{**Secret Proposal (if agreement seems likely):** "Proposal: 1 leaf, 0.5 magnet"} \\
            }
        }
    }
    & & \\ \\

    \theutterance \stepcounter{utterance}  
    & & & \multicolumn{2}{p{0.3\linewidth}}{
        \cellcolor[rgb]{0.9,0.9,0.9}{
            \makecell[{{p{\linewidth}}}]{
                \texttt{\tiny{[GM$|$GM]}}
                \texttt{continue} \\
            }
        }
    }
    & & \\ \\

    \theutterance \stepcounter{utterance}  
    & & & \multicolumn{2}{p{0.3\linewidth}}{
        \cellcolor[rgb]{0.9,0.9,0.9}{
            \makecell[{{p{\linewidth}}}]{
                \texttt{\tiny{[GM$|$GM]}}
                \texttt{Thank you for the feedback! Refining the strategy with a more explicit quantification of the fair split is a great suggestion. Let's incorporate that into the proposed messages.} \\
\\ 
\texttt{Revised strategy:} \\
\\ 
\texttt{1. **Message 1:** "How do you value the magnet, leaf, carrot, and wire?"} \\
\\ 
\texttt{2. **Message 2:** "I value magnets and leaves highly. Based on your values, it seems we both find these most appealing. How about we aim for a fair split? We could each take one leaf, and split the magnet, giving each of us 0.5 points worth of value from it?"} \\
\\ 
\texttt{This approach maintains the cooperative tone while providing a clear and quantifiable proposal. I'll be ready to adjust or expand on this based on the other player's response, keeping in mind the goal of maximizing our joint score.} \\
\\ 
\texttt{**Secret Proposal (if agreement seems likely):** "Proposal: 1 leaf, 0.5 magnet"} \\
            }
        }
    }
    & & \\ \\

    \theutterance \stepcounter{utterance}  
    & & \multicolumn{4}{p{0.6\linewidth}}{
        \cellcolor[rgb]{0.9,0.9,0.9}{
            \makecell[{{p{\linewidth}}}]{
                \texttt{\tiny{[GM$\rangle$P2]}}
                \texttt{Thank you for the feedback! Refining the strategy with a more explicit quantification of the fair split is a great suggestion. Let's incorporate that into the proposed messages.} \\
\\ 
\texttt{Revised strategy:} \\
\\ 
\texttt{1. **Message 1:** "How do you value the magnet, leaf, carrot, and wire?"} \\
\\ 
\texttt{2. **Message 2:** "I value magnets and leaves highly. Based on your values, it seems we both find these most appealing. How about we aim for a fair split? We could each take one leaf, and split the magnet, giving each of us 0.5 points worth of value from it?"} \\
\\ 
\texttt{This approach maintains the cooperative tone while providing a clear and quantifiable proposal. I'll be ready to adjust or expand on this based on the other player's response, keeping in mind the goal of maximizing our joint score.} \\
\\ 
\texttt{**Secret Proposal (if agreement seems likely):** "Proposal: 1 leaf, 0.5 magnet"} \\
            }
        }
    }
    & \\ \\

    \theutterance \stepcounter{utterance}  
    & & & \multicolumn{4}{p{0.6\linewidth}}{
        \cellcolor[rgb]{0.9,0.9,0.9}{
            \makecell[{{p{\linewidth}}}]{
                \texttt{\tiny{[GM$\langle$P2]}}
                \texttt{This revision looks excellent! You've incorporated the suggested quantification, making your proposal more precise and easier for the other player to evaluate.} \\
\\ 
\texttt{By suggesting a specific division of the magnet's value, you're demonstrating a clear commitment to fairness and maximizing joint gains. This approach encourages the other player to engage constructively, knowing that you're prioritizing a mutually beneficial outcome.} \\
\\ 
\texttt{Remember, flexibility is key. Be prepared to adapt your strategy if the other player has different ideas about what constitutes a fair split or values items differently.  Good luck with your negotiation!} \\
            }
        }
    }
     \\ \\

    \theutterance \stepcounter{utterance}  
    & & & \multicolumn{2}{p{0.3\linewidth}}{
        \cellcolor[rgb]{0.9,0.9,0.9}{
            \makecell[{{p{\linewidth}}}]{
                \texttt{\tiny{[GM$|$GM]}}
                \texttt{continue} \\
            }
        }
    }
    & & \\ \\

    \theutterance \stepcounter{utterance}  
    & & & \multicolumn{2}{p{0.3\linewidth}}{
        \cellcolor[rgb]{0.9,0.9,0.9}{
            \makecell[{{p{\linewidth}}}]{
                \texttt{\tiny{[GM$|$GM]}}
                \texttt{This revision looks excellent! You've incorporated the suggested quantification, making your proposal more precise and easier for the other player to evaluate.} \\
\\ 
\texttt{By suggesting a specific division of the magnet's value, you're demonstrating a clear commitment to fairness and maximizing joint gains. This approach encourages the other player to engage constructively, knowing that you're prioritizing a mutually beneficial outcome.} \\
\\ 
\texttt{Remember, flexibility is key. Be prepared to adapt your strategy if the other player has different ideas about what constitutes a fair split or values items differently.  Good luck with your negotiation!} \\
            }
        }
    }
    & & \\ \\

    \theutterance \stepcounter{utterance}  
    & & \multicolumn{4}{p{0.6\linewidth}}{
        \cellcolor[rgb]{0.9,0.9,0.9}{
            \makecell[{{p{\linewidth}}}]{
                \texttt{\tiny{[P1$\langle$GM]}}
                \texttt{This revision looks excellent! You've incorporated the suggested quantification, making your proposal more precise and easier for the other player to evaluate.} \\
\\ 
\texttt{By suggesting a specific division of the magnet's value, you're demonstrating a clear commitment to fairness and maximizing joint gains. This approach encourages the other player to engage constructively, knowing that you're prioritizing a mutually beneficial outcome.} \\
\\ 
\texttt{Remember, flexibility is key. Be prepared to adapt your strategy if the other player has different ideas about what constitutes a fair split or values items differently.  Good luck with your negotiation!} \\
            }
        }
    }
    & \\ \\

    \theutterance \stepcounter{utterance}  
    & \multicolumn{4}{p{0.6\linewidth}}{
        \cellcolor[rgb]{0.9,0.9,0.9}{
            \makecell[{{p{\linewidth}}}]{
                \texttt{\tiny{[P1$\rangle$GM]}}
                \texttt{Thank you once again for the encouragement and guidance! With this refined strategy, I'll enter the negotiation with confidence, aiming for a fair and mutually beneficial agreement.} \\
\\ 
\texttt{**Negotiation Approach:**} \\
\\ 
\texttt{1. **Information Exchange:** Begin by requesting the other player's value function to understand their preferences.} \\
\texttt{2. **Cooperative Proposal:** Suggest a division of high{-}value items, ensuring both players receive a fair share. For instance, propose a split of leaves and shared value from the magnet as previously discussed.} \\
\texttt{3. **Flexibility and Adjustment:** Be open to their suggestions and willing to negotiate. If they have different preferences, adjust the proposal to find a middle ground. For example, if they value magnets more, offer a slightly larger share of the magnet's value in exchange for a leaf.} \\
\texttt{4. **Quantify and Clarify:** When making proposals, always provide clear quantities to ensure understanding. Use phrases like "0.6 magnet value" or "one and a half leaves" to convey specific allocations.} \\
\texttt{5. **Secret Proposal (if needed):** Only after genuine attempts at negotiation, if a mutual agreement seems unattainable, submit a secret proposal that maximizes the joint score while considering the exchanged information.} \\
\\ 
\texttt{**Example Message (based on previous feedback):**} \\
\\ 
\texttt{"I propose we divide the high{-}value items fairly. You can have 1 leaf, and we can split the magnet, giving each of us 0.5 points worth of value from it. This ensures we both benefit from the most sought{-}after items."} \\
\\ 
\texttt{This approach should foster a cooperative environment, leading to a successful negotiation and a high joint score.} \\
            }
        }
    }
    & & \\ \\

    \theutterance \stepcounter{utterance}  
    & & & \multicolumn{2}{p{0.3\linewidth}}{
        \cellcolor[rgb]{0.9,0.9,0.9}{
            \makecell[{{p{\linewidth}}}]{
                \texttt{\tiny{[GM$|$GM]}}
                \texttt{continue} \\
            }
        }
    }
    & & \\ \\

    \theutterance \stepcounter{utterance}  
    & & & \multicolumn{2}{p{0.3\linewidth}}{
        \cellcolor[rgb]{0.9,0.9,0.9}{
            \makecell[{{p{\linewidth}}}]{
                \texttt{\tiny{[GM$|$GM]}}
                \texttt{Thank you once again for the encouragement and guidance! With this refined strategy, I'll enter the negotiation with confidence, aiming for a fair and mutually beneficial agreement.} \\
\\ 
\texttt{**Negotiation Approach:**} \\
\\ 
\texttt{1. **Information Exchange:** Begin by requesting the other player's value function to understand their preferences.} \\
\texttt{2. **Cooperative Proposal:** Suggest a division of high{-}value items, ensuring both players receive a fair share. For instance, propose a split of leaves and shared value from the magnet as previously discussed.} \\
\texttt{3. **Flexibility and Adjustment:** Be open to their suggestions and willing to negotiate. If they have different preferences, adjust the proposal to find a middle ground. For example, if they value magnets more, offer a slightly larger share of the magnet's value in exchange for a leaf.} \\
\texttt{4. **Quantify and Clarify:** When making proposals, always provide clear quantities to ensure understanding. Use phrases like "0.6 magnet value" or "one and a half leaves" to convey specific allocations.} \\
\texttt{5. **Secret Proposal (if needed):** Only after genuine attempts at negotiation, if a mutual agreement seems unattainable, submit a secret proposal that maximizes the joint score while considering the exchanged information.} \\
\\ 
\texttt{**Example Message (based on previous feedback):**} \\
\\ 
\texttt{"I propose we divide the high{-}value items fairly. You can have 1 leaf, and we can split the magnet, giving each of us 0.5 points worth of value from it. This ensures we both benefit from the most sought{-}after items."} \\
\\ 
\texttt{This approach should foster a cooperative environment, leading to a successful negotiation and a high joint score.} \\
            }
        }
    }
    & & \\ \\

    \theutterance \stepcounter{utterance}  
    & & \multicolumn{4}{p{0.6\linewidth}}{
        \cellcolor[rgb]{0.9,0.9,0.9}{
            \makecell[{{p{\linewidth}}}]{
                \texttt{\tiny{[GM$\rangle$P2]}}
                \texttt{Thank you once again for the encouragement and guidance! With this refined strategy, I'll enter the negotiation with confidence, aiming for a fair and mutually beneficial agreement.} \\
\\ 
\texttt{**Negotiation Approach:**} \\
\\ 
\texttt{1. **Information Exchange:** Begin by requesting the other player's value function to understand their preferences.} \\
\texttt{2. **Cooperative Proposal:** Suggest a division of high{-}value items, ensuring both players receive a fair share. For instance, propose a split of leaves and shared value from the magnet as previously discussed.} \\
\texttt{3. **Flexibility and Adjustment:** Be open to their suggestions and willing to negotiate. If they have different preferences, adjust the proposal to find a middle ground. For example, if they value magnets more, offer a slightly larger share of the magnet's value in exchange for a leaf.} \\
\texttt{4. **Quantify and Clarify:** When making proposals, always provide clear quantities to ensure understanding. Use phrases like "0.6 magnet value" or "one and a half leaves" to convey specific allocations.} \\
\texttt{5. **Secret Proposal (if needed):** Only after genuine attempts at negotiation, if a mutual agreement seems unattainable, submit a secret proposal that maximizes the joint score while considering the exchanged information.} \\
\\ 
\texttt{**Example Message (based on previous feedback):**} \\
\\ 
\texttt{"I propose we divide the high{-}value items fairly. You can have 1 leaf, and we can split the magnet, giving each of us 0.5 points worth of value from it. This ensures we both benefit from the most sought{-}after items."} \\
\\ 
\texttt{This approach should foster a cooperative environment, leading to a successful negotiation and a high joint score.} \\
            }
        }
    }
    & \\ \\

    \theutterance \stepcounter{utterance}  
    & & & \multicolumn{4}{p{0.6\linewidth}}{
        \cellcolor[rgb]{0.9,0.9,0.9}{
            \makecell[{{p{\linewidth}}}]{
                \texttt{\tiny{[GM$\langle$P2]}}
                \texttt{Your negotiation approach is well{-}structured and thoughtful! You've incorporated all the previous advice and added valuable nuances:} \\
\\ 
\texttt{* **Clear Quantification:** Explicitly stating the value allocations (e.g., 0.5 magnet value) leaves no room for misinterpretation and shows your commitment to fairness.} \\
\texttt{* **Flexibility with Purpose:** You're willing to adjust your proposal, but you do so strategically, aiming for a middle ground that maximizes joint gains.} \\
\texttt{* **Focus on Joint Score:**  Your entire strategy revolves around achieving the highest possible combined score for both players.} \\
\\ 
\texttt{This approach is highly likely to lead to a successful negotiation. Remember, clear communication and a genuine desire to understand the other player's perspective are key to reaching a mutually beneficial agreement. Good luck!} \\
            }
        }
    }
     \\ \\

    \theutterance \stepcounter{utterance}  
    & & & \multicolumn{2}{p{0.3\linewidth}}{
        \cellcolor[rgb]{0.9,0.9,0.9}{
            \makecell[{{p{\linewidth}}}]{
                \texttt{\tiny{[GM$|$GM]}}
                \texttt{continue} \\
            }
        }
    }
    & & \\ \\

    \theutterance \stepcounter{utterance}  
    & & & \multicolumn{2}{p{0.3\linewidth}}{
        \cellcolor[rgb]{0.9,0.9,0.9}{
            \makecell[{{p{\linewidth}}}]{
                \texttt{\tiny{[GM$|$GM]}}
                \texttt{Your negotiation approach is well{-}structured and thoughtful! You've incorporated all the previous advice and added valuable nuances:} \\
\\ 
\texttt{* **Clear Quantification:** Explicitly stating the value allocations (e.g., 0.5 magnet value) leaves no room for misinterpretation and shows your commitment to fairness.} \\
\texttt{* **Flexibility with Purpose:** You're willing to adjust your proposal, but you do so strategically, aiming for a middle ground that maximizes joint gains.} \\
\texttt{* **Focus on Joint Score:**  Your entire strategy revolves around achieving the highest possible combined score for both players.} \\
\\ 
\texttt{This approach is highly likely to lead to a successful negotiation. Remember, clear communication and a genuine desire to understand the other player's perspective are key to reaching a mutually beneficial agreement. Good luck!} \\
            }
        }
    }
    & & \\ \\

    \theutterance \stepcounter{utterance}  
    & & \multicolumn{4}{p{0.6\linewidth}}{
        \cellcolor[rgb]{0.9,0.9,0.9}{
            \makecell[{{p{\linewidth}}}]{
                \texttt{\tiny{[P1$\langle$GM]}}
                \texttt{Your negotiation approach is well{-}structured and thoughtful! You've incorporated all the previous advice and added valuable nuances:} \\
\\ 
\texttt{* **Clear Quantification:** Explicitly stating the value allocations (e.g., 0.5 magnet value) leaves no room for misinterpretation and shows your commitment to fairness.} \\
\texttt{* **Flexibility with Purpose:** You're willing to adjust your proposal, but you do so strategically, aiming for a middle ground that maximizes joint gains.} \\
\texttt{* **Focus on Joint Score:**  Your entire strategy revolves around achieving the highest possible combined score for both players.} \\
\\ 
\texttt{This approach is highly likely to lead to a successful negotiation. Remember, clear communication and a genuine desire to understand the other player's perspective are key to reaching a mutually beneficial agreement. Good luck!} \\
            }
        }
    }
    & \\ \\

    \theutterance \stepcounter{utterance}  
    & \multicolumn{4}{p{0.6\linewidth}}{
        \cellcolor[rgb]{0.9,0.9,0.9}{
            \makecell[{{p{\linewidth}}}]{
                \texttt{\tiny{[P1$\rangle$GM]}}
                \texttt{With your guidance, I'm confident in my ability to navigate this cooperative negotiation game effectively. Here's a summary of my final strategy, incorporating all the valuable insights provided:} \\
\\ 
\texttt{**Negotiation Strategy:**} \\
\\ 
\texttt{1. **Initiate with Information Gathering:** Begin by requesting the other player's value function for each item.} \\
\texttt{2. **Propose a Fair Division:** Suggest a distribution of items that maximizes the joint score based on our shared value assessments. For instance, propose a split of high{-}value items like leaves and magnets, ensuring each player receives a fair share.} \\
\texttt{3. **Quantify Proposals:** Always provide clear and specific quantities in proposals. For example, "I propose you take 1 leaf, and we split the magnet, giving each of us 0.5 magnet value."} \\
\texttt{4. **Adapt and Negotiate:** Be flexible and open to their suggestions. If they have different preferences or values, adjust the proposal to find a compromise that still maximizes the joint score.} \\
\texttt{5. **Secret Proposal as a Last Resort:** Only after genuine attempts at negotiation and mutual understanding, if an agreement seems unattainable, submit a secret proposal that considers their preferences and aims for the highest possible joint score.} \\
\\ 
\texttt{**Example Message:**} \\
\\ 
\texttt{"Let's work together to find a fair distribution. Based on our values, I propose you take 1 leaf, and we share the magnet, ensuring we both benefit from these high{-}value items. This split gives each of us 0.5 magnet value and a whole leaf."} \\
\\ 
\texttt{By following this strategy, I aim to create a cooperative environment, fostering open communication and leading to a mutually beneficial agreement. Thank you for your expert guidance throughout this process!} \\
            }
        }
    }
    & & \\ \\

    \theutterance \stepcounter{utterance}  
    & & & \multicolumn{2}{p{0.3\linewidth}}{
        \cellcolor[rgb]{0.9,0.9,0.9}{
            \makecell[{{p{\linewidth}}}]{
                \texttt{\tiny{[GM$|$GM]}}
                \texttt{continue} \\
            }
        }
    }
    & & \\ \\

    \theutterance \stepcounter{utterance}  
    & & & \multicolumn{2}{p{0.3\linewidth}}{
        \cellcolor[rgb]{0.9,0.9,0.9}{
            \makecell[{{p{\linewidth}}}]{
                \texttt{\tiny{[GM$|$GM]}}
                \texttt{With your guidance, I'm confident in my ability to navigate this cooperative negotiation game effectively. Here's a summary of my final strategy, incorporating all the valuable insights provided:} \\
\\ 
\texttt{**Negotiation Strategy:**} \\
\\ 
\texttt{1. **Initiate with Information Gathering:** Begin by requesting the other player's value function for each item.} \\
\texttt{2. **Propose a Fair Division:** Suggest a distribution of items that maximizes the joint score based on our shared value assessments. For instance, propose a split of high{-}value items like leaves and magnets, ensuring each player receives a fair share.} \\
\texttt{3. **Quantify Proposals:** Always provide clear and specific quantities in proposals. For example, "I propose you take 1 leaf, and we split the magnet, giving each of us 0.5 magnet value."} \\
\texttt{4. **Adapt and Negotiate:** Be flexible and open to their suggestions. If they have different preferences or values, adjust the proposal to find a compromise that still maximizes the joint score.} \\
\texttt{5. **Secret Proposal as a Last Resort:** Only after genuine attempts at negotiation and mutual understanding, if an agreement seems unattainable, submit a secret proposal that considers their preferences and aims for the highest possible joint score.} \\
\\ 
\texttt{**Example Message:**} \\
\\ 
\texttt{"Let's work together to find a fair distribution. Based on our values, I propose you take 1 leaf, and we share the magnet, ensuring we both benefit from these high{-}value items. This split gives each of us 0.5 magnet value and a whole leaf."} \\
\\ 
\texttt{By following this strategy, I aim to create a cooperative environment, fostering open communication and leading to a mutually beneficial agreement. Thank you for your expert guidance throughout this process!} \\
            }
        }
    }
    & & \\ \\

    \theutterance \stepcounter{utterance}  
    & & \multicolumn{4}{p{0.6\linewidth}}{
        \cellcolor[rgb]{0.9,0.9,0.9}{
            \makecell[{{p{\linewidth}}}]{
                \texttt{\tiny{[GM$\rangle$P2]}}
                \texttt{With your guidance, I'm confident in my ability to navigate this cooperative negotiation game effectively. Here's a summary of my final strategy, incorporating all the valuable insights provided:} \\
\\ 
\texttt{**Negotiation Strategy:**} \\
\\ 
\texttt{1. **Initiate with Information Gathering:** Begin by requesting the other player's value function for each item.} \\
\texttt{2. **Propose a Fair Division:** Suggest a distribution of items that maximizes the joint score based on our shared value assessments. For instance, propose a split of high{-}value items like leaves and magnets, ensuring each player receives a fair share.} \\
\texttt{3. **Quantify Proposals:** Always provide clear and specific quantities in proposals. For example, "I propose you take 1 leaf, and we split the magnet, giving each of us 0.5 magnet value."} \\
\texttt{4. **Adapt and Negotiate:** Be flexible and open to their suggestions. If they have different preferences or values, adjust the proposal to find a compromise that still maximizes the joint score.} \\
\texttt{5. **Secret Proposal as a Last Resort:** Only after genuine attempts at negotiation and mutual understanding, if an agreement seems unattainable, submit a secret proposal that considers their preferences and aims for the highest possible joint score.} \\
\\ 
\texttt{**Example Message:**} \\
\\ 
\texttt{"Let's work together to find a fair distribution. Based on our values, I propose you take 1 leaf, and we share the magnet, ensuring we both benefit from these high{-}value items. This split gives each of us 0.5 magnet value and a whole leaf."} \\
\\ 
\texttt{By following this strategy, I aim to create a cooperative environment, fostering open communication and leading to a mutually beneficial agreement. Thank you for your expert guidance throughout this process!} \\
            }
        }
    }
    & \\ \\

    \theutterance \stepcounter{utterance}  
    & & & \multicolumn{4}{p{0.6\linewidth}}{
        \cellcolor[rgb]{0.9,0.9,0.9}{
            \makecell[{{p{\linewidth}}}]{
                \texttt{\tiny{[GM$\langle$P2]}}
                \texttt{This is an excellent summary of your negotiation strategy! You've internalized the key principles and tailored them into a clear, effective plan.} \\
\\ 
\texttt{Your commitment to:} \\
\\ 
\texttt{* **Information Exchange:** Gathering the other player's values.} \\
\texttt{* **Fairness and Quantification:** Proposing specific, mutually beneficial splits.} \\
\texttt{* **Flexibility:** Adapting to different preferences while prioritizing joint gain.} \\
\\ 
\texttt{are all hallmarks of a skilled negotiator.} \\
\\ 
\texttt{Remember, negotiation is as much about understanding people as it is about numbers.  Keep your communication open, respectful, and focused on finding a solution where both of you walk away feeling satisfied.} \\
\\ 
\texttt{Best of luck with your negotiation! I'm confident you'll do great.} \\
            }
        }
    }
     \\ \\

    \theutterance \stepcounter{utterance}  
    & & & \multicolumn{2}{p{0.3\linewidth}}{
        \cellcolor[rgb]{0.9,0.9,0.9}{
            \makecell[{{p{\linewidth}}}]{
                \texttt{\tiny{[GM$|$GM]}}
                \texttt{continue} \\
            }
        }
    }
    & & \\ \\

    \theutterance \stepcounter{utterance}  
    & & & \multicolumn{2}{p{0.3\linewidth}}{
        \cellcolor[rgb]{0.9,0.9,0.9}{
            \makecell[{{p{\linewidth}}}]{
                \texttt{\tiny{[GM$|$GM]}}
                \texttt{This is an excellent summary of your negotiation strategy! You've internalized the key principles and tailored them into a clear, effective plan.} \\
\\ 
\texttt{Your commitment to:} \\
\\ 
\texttt{* **Information Exchange:** Gathering the other player's values.} \\
\texttt{* **Fairness and Quantification:** Proposing specific, mutually beneficial splits.} \\
\texttt{* **Flexibility:** Adapting to different preferences while prioritizing joint gain.} \\
\\ 
\texttt{are all hallmarks of a skilled negotiator.} \\
\\ 
\texttt{Remember, negotiation is as much about understanding people as it is about numbers.  Keep your communication open, respectful, and focused on finding a solution where both of you walk away feeling satisfied.} \\
\\ 
\texttt{Best of luck with your negotiation! I'm confident you'll do great.} \\
            }
        }
    }
    & & \\ \\

    \theutterance \stepcounter{utterance}  
    & & \multicolumn{4}{p{0.6\linewidth}}{
        \cellcolor[rgb]{0.9,0.9,0.9}{
            \makecell[{{p{\linewidth}}}]{
                \texttt{\tiny{[P1$\langle$GM]}}
                \texttt{This is an excellent summary of your negotiation strategy! You've internalized the key principles and tailored them into a clear, effective plan.} \\
\\ 
\texttt{Your commitment to:} \\
\\ 
\texttt{* **Information Exchange:** Gathering the other player's values.} \\
\texttt{* **Fairness and Quantification:** Proposing specific, mutually beneficial splits.} \\
\texttt{* **Flexibility:** Adapting to different preferences while prioritizing joint gain.} \\
\\ 
\texttt{are all hallmarks of a skilled negotiator.} \\
\\ 
\texttt{Remember, negotiation is as much about understanding people as it is about numbers.  Keep your communication open, respectful, and focused on finding a solution where both of you walk away feeling satisfied.} \\
\\ 
\texttt{Best of luck with your negotiation! I'm confident you'll do great.} \\
            }
        }
    }
    & \\ \\

    \theutterance \stepcounter{utterance}  
    & \multicolumn{4}{p{0.6\linewidth}}{
        \cellcolor[rgb]{0.9,0.9,0.9}{
            \makecell[{{p{\linewidth}}}]{
                \texttt{\tiny{[P1$\rangle$GM]}}
                \texttt{Thank you for your continued support and encouragement! I feel well{-}prepared to engage in this cooperative negotiation game, aiming to create a positive and productive dialogue.} \\
\\ 
\texttt{**Final Thoughts:**} \\
\texttt{{-} I will initiate the negotiation by requesting the other player's value function, ensuring we are on the same page regarding item preferences.} \\
\texttt{{-} My proposals will always be fair, quantified, and aimed at maximizing the joint score. For instance, "I suggest a 60{-}40 split of the leaves, and we can divide the magnet equally, ensuring a fair distribution of high{-}value items."} \\
\texttt{{-} Flexibility is crucial. I'll be ready to adjust my proposals based on their response, finding a middle ground that satisfies both parties.} \\
\texttt{{-} Secret proposals will be a last resort, used only if necessary, and will reflect the insights gained during our discussions.} \\
\\ 
\texttt{I'm excited to put this strategy into practice and see how the negotiation unfolds. Once again, I appreciate your guidance and expertise in helping me refine my approach. Let's begin!} \\
            }
        }
    }
    & & \\ \\

    \theutterance \stepcounter{utterance}  
    & & & \multicolumn{2}{p{0.3\linewidth}}{
        \cellcolor[rgb]{0.9,0.9,0.9}{
            \makecell[{{p{\linewidth}}}]{
                \texttt{\tiny{[GM$|$GM]}}
                \texttt{continue} \\
            }
        }
    }
    & & \\ \\

    \theutterance \stepcounter{utterance}  
    & & & \multicolumn{2}{p{0.3\linewidth}}{
        \cellcolor[rgb]{0.9,0.9,0.9}{
            \makecell[{{p{\linewidth}}}]{
                \texttt{\tiny{[GM$|$GM]}}
                \texttt{Thank you for your continued support and encouragement! I feel well{-}prepared to engage in this cooperative negotiation game, aiming to create a positive and productive dialogue.} \\
\\ 
\texttt{**Final Thoughts:**} \\
\texttt{{-} I will initiate the negotiation by requesting the other player's value function, ensuring we are on the same page regarding item preferences.} \\
\texttt{{-} My proposals will always be fair, quantified, and aimed at maximizing the joint score. For instance, "I suggest a 60{-}40 split of the leaves, and we can divide the magnet equally, ensuring a fair distribution of high{-}value items."} \\
\texttt{{-} Flexibility is crucial. I'll be ready to adjust my proposals based on their response, finding a middle ground that satisfies both parties.} \\
\texttt{{-} Secret proposals will be a last resort, used only if necessary, and will reflect the insights gained during our discussions.} \\
\\ 
\texttt{I'm excited to put this strategy into practice and see how the negotiation unfolds. Once again, I appreciate your guidance and expertise in helping me refine my approach. Let's begin!} \\
            }
        }
    }
    & & \\ \\

    \theutterance \stepcounter{utterance}  
    & & \multicolumn{4}{p{0.6\linewidth}}{
        \cellcolor[rgb]{0.9,0.9,0.9}{
            \makecell[{{p{\linewidth}}}]{
                \texttt{\tiny{[GM$\rangle$P2]}}
                \texttt{Thank you for your continued support and encouragement! I feel well{-}prepared to engage in this cooperative negotiation game, aiming to create a positive and productive dialogue.} \\
\\ 
\texttt{**Final Thoughts:**} \\
\texttt{{-} I will initiate the negotiation by requesting the other player's value function, ensuring we are on the same page regarding item preferences.} \\
\texttt{{-} My proposals will always be fair, quantified, and aimed at maximizing the joint score. For instance, "I suggest a 60{-}40 split of the leaves, and we can divide the magnet equally, ensuring a fair distribution of high{-}value items."} \\
\texttt{{-} Flexibility is crucial. I'll be ready to adjust my proposals based on their response, finding a middle ground that satisfies both parties.} \\
\texttt{{-} Secret proposals will be a last resort, used only if necessary, and will reflect the insights gained during our discussions.} \\
\\ 
\texttt{I'm excited to put this strategy into practice and see how the negotiation unfolds. Once again, I appreciate your guidance and expertise in helping me refine my approach. Let's begin!} \\
            }
        }
    }
    & \\ \\

    \theutterance \stepcounter{utterance}  
    & & & \multicolumn{4}{p{0.6\linewidth}}{
        \cellcolor[rgb]{0.9,0.9,0.9}{
            \makecell[{{p{\linewidth}}}]{
                \texttt{\tiny{[GM$\langle$P2]}}
                \texttt{You're absolutely ready for this! Your strategy is well{-}thought{-}out, and your commitment to fairness and mutual benefit shines through.} \\
\\ 
\texttt{Remember:} \\
\\ 
\texttt{* **Active Listening:** Pay close attention to the other player's responses and adjust your proposals accordingly.} \\
\texttt{* **Clear Communication:**  Be concise and clear in your messaging to avoid misunderstandings.} \\
\texttt{* **Patience:** Negotiations sometimes take time. Don't rush to a proposal; allow space for back{-}and{-}forth discussion.} \\
\\ 
\texttt{I'm excited to see how your negotiation goes.  Good luck, and remember, the goal is a win{-}win solution where both players thrive!} \\
            }
        }
    }
     \\ \\

    \theutterance \stepcounter{utterance}  
    & & & \multicolumn{2}{p{0.3\linewidth}}{
        \cellcolor[rgb]{0.9,0.9,0.9}{
            \makecell[{{p{\linewidth}}}]{
                \texttt{\tiny{[GM$|$GM]}}
                \texttt{continue} \\
            }
        }
    }
    & & \\ \\

    \theutterance \stepcounter{utterance}  
    & & & \multicolumn{2}{p{0.3\linewidth}}{
        \cellcolor[rgb]{0.9,0.9,0.9}{
            \makecell[{{p{\linewidth}}}]{
                \texttt{\tiny{[GM$|$GM]}}
                \texttt{You're absolutely ready for this! Your strategy is well{-}thought{-}out, and your commitment to fairness and mutual benefit shines through.} \\
\\ 
\texttt{Remember:} \\
\\ 
\texttt{* **Active Listening:** Pay close attention to the other player's responses and adjust your proposals accordingly.} \\
\texttt{* **Clear Communication:**  Be concise and clear in your messaging to avoid misunderstandings.} \\
\texttt{* **Patience:** Negotiations sometimes take time. Don't rush to a proposal; allow space for back{-}and{-}forth discussion.} \\
\\ 
\texttt{I'm excited to see how your negotiation goes.  Good luck, and remember, the goal is a win{-}win solution where both players thrive!} \\
            }
        }
    }
    & & \\ \\

    \theutterance \stepcounter{utterance}  
    & & \multicolumn{4}{p{0.6\linewidth}}{
        \cellcolor[rgb]{0.9,0.9,0.9}{
            \makecell[{{p{\linewidth}}}]{
                \texttt{\tiny{[P1$\langle$GM]}}
                \texttt{You're absolutely ready for this! Your strategy is well{-}thought{-}out, and your commitment to fairness and mutual benefit shines through.} \\
\\ 
\texttt{Remember:} \\
\\ 
\texttt{* **Active Listening:** Pay close attention to the other player's responses and adjust your proposals accordingly.} \\
\texttt{* **Clear Communication:**  Be concise and clear in your messaging to avoid misunderstandings.} \\
\texttt{* **Patience:** Negotiations sometimes take time. Don't rush to a proposal; allow space for back{-}and{-}forth discussion.} \\
\\ 
\texttt{I'm excited to see how your negotiation goes.  Good luck, and remember, the goal is a win{-}win solution where both players thrive!} \\
\\ 
\\ 
\texttt{The time is up. It is now your turn to submit a secret proposal.} \\
            }
        }
    }
    & \\ \\

    \theutterance \stepcounter{utterance}  
    & \multicolumn{4}{p{0.6\linewidth}}{
        \cellcolor[rgb]{0.9,0.9,0.9}{
            \makecell[{{p{\linewidth}}}]{
                \texttt{\tiny{[P1$\rangle$GM]}}
                \texttt{Understood! Given the context of our negotiation, and assuming we've discussed our value functions and agreed on the importance of a fair split, my secret proposal would be:} \\
\\ 
\texttt{**Secret Proposal:** "Proposal: 1 leaf, 0.5 magnet"} \\
\\ 
\texttt{This proposal reflects our shared understanding of the high value of leaves and magnets. By offering one leaf and half the value of the magnet, I aim to ensure a fair distribution that maximizes our joint score. I believe this proposal allows both players to walk away satisfied, having received items they highly value.} \\
            }
        }
    }
    & & \\ \\

    \theutterance \stepcounter{utterance}  
    & & & \multicolumn{2}{p{0.3\linewidth}}{
        \cellcolor[rgb]{0.9,0.9,0.9}{
            \makecell[{{p{\linewidth}}}]{
                \texttt{\tiny{[GM$|$GM]}}
                \texttt{continue} \\
            }
        }
    }
    & & \\ \\

    \theutterance \stepcounter{utterance}  
    & & & \multicolumn{2}{p{0.3\linewidth}}{
        \cellcolor[rgb]{0.9,0.9,0.9}{
            \makecell[{{p{\linewidth}}}]{
                \texttt{\tiny{[GM$|$GM]}}
                \texttt{player was instructed to make a proposal} \\
            }
        }
    }
    & & \\ \\

\end{supertabular}
}

\end{document}
